\DTMsavedate{mydate}{2022-06-30}
\maketitle{Embedded Systems}{Electronic Engineering}{\DTMusedate{mydate}}

% ----------------------------------------------------- %
% COURSE INFO
\subsection*{Course Information}\label{course:3232}
\vspace{0ex}%
\noindent\parbox{0.5\textwidth}{%
    \noindent\begin{tabular}{@{}ll}
                 \textsf{Course:}          & 	Embedded Systems  \\
                 \textsf{Course Code:}         & 	TEE 3232    \\
                 \textsf{Credits:}    & 	10
    \end{tabular}}
\vspace{2ex}\hrule\vspace{2ex}

% ----------------------------------------------------- %
% INSTRUCTOR INFO
\subsection*{Lecturer Information}
\noindent\parbox{0.5\textwidth}{%
    \noindent\begin{tabular}{@{}ll}
                 \textsf{Name:}         &     Reginald Gonye         \\
%\textsf{Phone:}        &    \phone          \\
\textsf{Email:}        &    reginald.gonye@nust.ac.zw \\
\textsf{Qualifications:}        &  MPhil in Electronic Engineering, (Hon) BEng in Electronic Engineering
    \end{tabular}}
\vspace{2ex}\hrule\vspace{2ex}


% ----------------------------------------------------- %
% COURSE SYNOPSIS
\subsection*{Course Synopsis}
The module explores applications of embedded systems in Microcontrollers i.e\. its memory maps, programming languages, I/O, timers, interrupts, hardware interfacing.
It also covers Picocontrollers i.e.\ memory maps, SFRs, stacks, programming languages, oscillator types, configuration fuses, watchdog timers and code protection.
\vspace{2ex}\hrule\vspace{2ex}

% ----------------------------------------------------- %
% LEARNING AIM
\subsection*{Learning Aim}
To instill concepts of Embedded Computer Systems into students.
\vspace{2ex} \hrule\vspace{2ex}

% ----------------------------------------------------- %
% LEARNING OBJECTIVES
\subsection*{Learning Objectives}
By the end of the course, candidates should be able to:
\begin{outline}
    \1      Understand the concept of embedded computers.

    \1  	Design an embedded computer solution for a given situation.

\end{outline}
\vspace{2ex}\hrule\vspace{2ex}

% ----------------------------------------------------- %
% COURSE TOPICS
\subsection*{Topics}
\begin{outline}
    \1 Introduction
    \1 The 8051 Microcontroller
    \begin{enumerate}
        \item Introduction
        \item Memory of the 8051
        \item Programming the 8051
        \item IO of the 8051
        \item Counters and timers
        \item Interrupts
    \end{enumerate}
    \1 The PIC15F84
    \begin{enumerate}
        \item Introduction
        \item Memory of the PIC15F84
        \item Programming the PIC15F84
        \item IO of the PIC15F84
        \item Counters and timers
        \item Interrupts
    \end{enumerate}
    \1 Root Locus
\end{outline}
\vspace{2ex} \hrule\vspace{2ex}

% ----------------------------------------------------- %
% Students Assenssment
\subsection*{Assessment of Students}
\begin{outline}
    \1 Continuous assessment will consist of tests, laboratory work and assignments (NB: Attendance may also contribute to the coursework mark).
    \1 Final assessment will consist of an open book examination at the end of the semester. The examination shall comprise of two sections: section A with three questions on microcontrollers and section B with three questions on PICs.  Each question shall carry 33.3 marks. You will be required to answer a total of three questions with at least one from each section. The questions will test your ability to design an embedded computer solution for a given situation.
\end{outline}
\vspace{2ex}\hrule\vspace{2ex}

% ----------------------------------------------------- %
% Grade Evaluation
\subsection*{Course Evaluation and Grading Policies}
Grades are based on:
Continuous assessment (\qty{25}{\percent}),
Final Exam (\qty{75}{\percent})
\vspace{2ex}\hrule\vspace{2ex}
% ----------------------------------------------------- %
% EQUIPMENT
% https://sites.udel.edu/computing-purchases/personal-specs/
% \subsection*{Essential Equipment}

% \vspace{2ex}\hrule\vspace{2ex}

% Policies and Other information
\subsection*{Policies and Other Information}
% Key Text or some Books
\subsection*{Key Text}

\begin{itemize}
    \item Lecture notes
\end{itemize}
\vspace{2ex}\hrule
\vspace{2ex}\hrule\vspace{2ex}